\phantomsection\bheading{section}{11.\quad Recommendations and open questions}{section}{\protect\numberline{11.}Recommendations and open questions}

\phantomsection\bheading{subsection}{11.1\quad Recommendation}{subsection}{\protect\numberline{11.1}Recommendation}

\begin{enumerate}
\item \textbf{Build the report mode} (§9) as a \code{check}-only pass with its dump and corpus fixtures, on the design of §2–§3, with the entry assumption P1–P6 for TEA programs read from the write-set pass. Nothing emitted changes. Budget: the effect pass's shape, a slice.
\item \textbf{Run it both ways} (\code{writers}, \code{rebuilds}) over \code{core/}, the corpus, TodoMVC, Conduit and the table app; classify by hand; report the table of §9.3 and the churn and timing rows.
\item \textbf{Decide the demand set first} (change 1). As asked, the rule touches nothing the owner's applications do; widened, it touches their every list edit. The number that decides is the rejection rate on \code{rebuilds}.
\item \textbf{Then decide error versus warning} on the exact/approximate split. The soundness argument holds for both; the choice is rule 7's.
\item \textbf{Write P1–P6 into \code{browser-direct.\allowbreak{}md}} before S7, whatever the decision: they are the obligations S7's in-place records need too.
\item If the error is adopted: specify the pass in \code{checker-v2.md} (an amendment, like §26), the \code{foreign} words in \code{boundary.md} §4, the demand summaries and \code{copy} in \code{language.md} §6.8, the representation change in \code{backend.md} §4, and the message catalogue in \code{language.md} §10 — then build it red-first, with the stability rule versioned in the interface format.
\end{enumerate}

\phantomsection\bheading{subsection}{11.2\quad What this document is confident of}{subsection}{\protect\numberline{11.2}What this document is confident of}

\begin{itemize}
\item The property (§2.3) is the right one: liveness-based, so reads before writes are free, and dead aliases are harmless. Every production system that accepted real array code has some form of it (Sisal's ordering edges, Clean's observing references, Futhark's syntactic order, Rust's two-phase borrows), and beni's strict order gives it with no construct.
\item The inference (§3) is directional, has a best solution, and is linear per round. The optimistic fixpoint with post-validation is the published recipe (AHK, Lean, Brandon).
\item The soundness argument (§4) has the right shape and the right assumptions; its lemmas are standard. The reviewer should attack A3/A4 (the platform) and §4.4's list.
\item Higher-order code is handled without annotation because lambdas are in front of the checker and function types already carry inferred classes (§3.6). This is the place every prior system gave up or declared; it is also where a reviewer should look hardest, at the $\mathit{supply}$/$\mathit{need}$ edges.
\end{itemize}

\phantomsection\bheading{subsection}{11.3\quad Open questions}{subsection}{\protect\numberline{11.3}Open questions}

\begin{enumerate}
\item \textbf{The rejection rate on real code}, exact and approximate, for both demand sets. Nothing in the literature stands in for it (research 69 §5.1). §9 answers it.
\item \textbf{The demand set} (change 1): the owner's call.
\item \textbf{\code{Ref} and \code{Queue}} (§5.12): a sound discipline for in-place edits of a cell's contents needs a whole-program argument about readers; left as a limit with a \code{copy} fix.
\item \textbf{Path-sensitive cell sets} (§5.6): worth their cost only if §9's \code{join} tag is frequent.
\item \textbf{Interface churn} under inferred summaries: measured by §9.3; change 5's optional annotation is the remedy if it is high.
\item \textbf{The warm budget}: a body edit that changes a summary is a signature edit; \code{plans/queue.md} row 55's signature-edit miss applies.
\item \textbf{Prepend without the trie} (change 8): decided by \code{prepends-loose}.
\item \textbf{The crash screen} (A9): confirm it prints no program value, or make it so.
\item \textbf{Record paths} (§7.7): the same rules, a stronger A4; S7's business.
\item \textbf{Konečný's best-typing proof} (AHK 2008 §6.2 cites Konečný 2003, not in the library) and the Cogent JFP 2021 paper (the formal “two semantics coincide” theorem O'Connor 2026 only states): both should be added to \code{references/\allowbreak{}ownership/} before the specification is written, so that §3.5's and §4.1's shapes can be checked against the originals.
\item \textbf{Exclusive contents in practice} (§2.9): how often the lists and dictionaries real programs edit through their elements are built exclusively, and which construction clears the flag most. If decoded data (every \code{List} a schema parses) is not declared element-fresh, nested edits of loaded data are all limits; the row for a decoder's output decides it.
\item \textbf{The two-slots-one-record hazard} is the one place this design relies on a flag the program's construction must establish rather than on liveness alone; a reviewer should try to break Lemma 1b with containers built through \code{where}-constrained generic code and through \code{foreign} siblings.
\item \textbf{Path-insensitive $\mathit{use}$} (§4.6's review): \code{f xs flag = if flag then List.\allowbreak{}length xs else List.\allowbreak{}length (List.\allowbreak{}push xs 1)} is $\mathsf{consumed}$, so \code{n = f xs True; List.\allowbreak{}length xs} is a false error. A conditional summary over a \emph{value} (not a class) is beyond §2.6's grammar; §9 counts how often it bites.
\item \textbf{The entry fixpoint's cost} (P5): one whole-program pass over \code{init} and the arms per build, the only non-modular computation here; its time goes beside the write-set pass's under \code{--self-profile}.
\item \textbf{A once-closure the program calls once in fact} passed to a $\mathsf{many}$ position stays an error; whether a \code{List.map} over a literal of one element deserves a special case is a question for the report mode's $\mathit{calls}\ \mathsf{many}$ tag count.
\end{enumerate}
